
\subsection{Graph domains of the Stokes linearization}
\label{r:sec:linear-domains}

The linearization around the smooth Stokes response has the
parabolic graph domain of the heat operator.
Write \(w=z\) for the prescribed Stokes response in
\eqref{r:joined-eq:stokeslift}, and define
\[
 K_wv=\Pp[(w\cdot\nabla)v+(v\cdot\nabla)w],
 \qquad L_w=\partial_t-\nu\Delta+K_w.
\]
For an integer \(m\ge0\), let
\[
 \mathcal E_m(0,S)=
 C([0,S];H^m_\sigma)\cap L^2(0,S;H^{m+1}_\sigma),
 \qquad v_t\in L^2(0,S;H^{m-1}_\sigma).
\]
The derivative condition is included in the definition of \(\mathcal E_m\).

\begin{proposition}[Linear graph domains]
For \(0<S\le T\), the initial-trace graph completions of \(L_w\) and
\(L_0=\partial_t-\nu\Delta\) both equal \(\mathcal E_m(0,S)\).
Their norms are equivalent with constants bounded independently of \(S\):
\begin{equation}
 \begin{aligned}
 &\norm v_{CH^m}+\norm v_{L^2H^{m+1}}
       +\norm{v_t}_{L^2H^{m-1}}\\
 &\qquad\le C_{m,\nu,T,w}
 \left(\norm{v(0)}_{H^m}+\norm{L_wv}_{L^2H^{m-1}}\right).
 \end{aligned}
 \label{r:eq:linear-domain-bound}
\end{equation}
Every coefficient in the constant is determined by finitely many prescribed
source norms, with no restriction on their size.
\end{proposition}

\begin{proof}
Product differentiation and the divergence form at \(m=0\) give
\[
 \norm{K_wv}_{H^{m-1}}
 \le C_m\norm w_{W^{m+1,\infty}}\norm v_{H^m}.
\]
Put \(r=L_wv\), and pair the equation at height \(m\), distributing
the duality between \(H^{m-1}\) and \(H^{m+1}\).
Since
\(\norm{\nabla v}_{H^m}^2=\norm v_{H^{m+1}}^2-\norm v_{H^m}^2\),
Young's inequality yields
\[
 \frac{\dd}{\dd t}\norm v_{H^m}^2+\nu\norm v_{H^{m+1}}^2
 \le C_{m,\nu,w}\norm v_{H^m}^2+C_\nu\norm r_{H^{m-1}}^2.
\]
Gronwall and integration give the first two norms in
\eqref{r:eq:linear-domain-bound}; the equation gives the third.
The exponential factor is bounded by its value at \(T\).
Conversely, each member of \(\mathcal E_m\) has \(L_wv\in L^2H^{m-1}\)
by the product bound. Approximation of its initial trace and residual,
followed by the difference estimate, gives the graph completion.
Existence for the approximating linear problems follows from the heat
integral equation on a finite partition, on which its smooth bounded
coefficients give contraction. The energy estimate joins those intervals.
\end{proof}

Temporal differentiation retains the earlier lift rows. If
\(W_n=\partial_t^nv\), \(Z_a=\partial_t^aw\), and
\(h_w=-\Pp[(w\cdot\nabla)w]\), the exact equation is
\[
 L_wW_n+\sum_{a=0}^n\binom na
       \Pp[(W_a\cdot\nabla)W_{n-a}]
 =\partial_t^nh_w-\sum_{a=1}^n\binom na K_{Z_a}W_{n-a}.
\]
For \(c_{a,m}=C_m\sup_{t\le T}\norm{Z_a}_{W^{m+1,\infty}}\),
\[
 \norm{\partial_t^n(K_wv)}_{L^2H^{m-1}}
 \le\sum_{a=0}^n\binom na c_{a,m}\norm{W_{n-a}}_{L^2H^m}.
\]
The adjacent pair has parabolic degree \(2n+m+1\); each velocity
factor on this right side has degree at most \(2n+m\).
At the base corner, kinetic energy already supplies the lower row.
These source contributions preserve the terminal linear domain.
The nonlinear self-interaction remains in the displayed equation, so
the domain estimate does not identify \(v\) with an unforced solution.
